\documentclass{article}
\usepackage[utf8]{inputenc}
\usepackage{amsmath,amssymb}
\usepackage[most]{tcolorbox}
\usepackage{geometry}
\geometry{margin=2.5cm}

\newcommand{\step}[1]{\par\noindent\hangindent=3.5em\hangafter=1%
  \makebox[3.5em][l]{\textbf{#1.}}}
\newcommand{\steptag}[1]{\hfill{\small\textsf{[#1]}}}

\newtcolorbox{proofbox}[1]{
  colback=gray!5, colframe=gray!60,
  fonttitle=\bfseries, title={#1},
  breakable, enhanced,
  left=4pt, right=4pt, top=4pt, bottom=4pt,
  before skip=10pt, after skip=10pt
}

\begin{document}

\begin{proofbox}{Polarised holes q\_\{r,s\}=P(r∘s)-r∘s obey, in the state...}
\step{1} Polarised holes q\_\{r,s\}=P(r∘s)-r∘s obey, in the state seminorm ||x||\_omega\textasciicircum{}2=omega(x\textasciicircum{}2), ||q\_\{r,s\}||\_omega <= C sqrt(eta)||r||||s|| (4.1) and ||q\_\{r,s\}|| <= C||r||||s||; hence ||P(q\_\{r,s\}∘q\_\{u,v\})|| <= C eta ||r||||s||||u||||v|| (HH) and ||P(q\_\{r,s\}∘z)|| <= C sqrt(eta)||r||||s||||z|| for z in V (HZ). \steptag{validated} \\[6pt]
\step{1.1} SETUP. Fix r,s,u,v in A=Im P (def-near-fixed-algebra) and z in V=B(H)\_sa. B(H)\_sa is the special Jordan algebra of self-adjoint operators and hence a JC algebra (GT-bhsa-jc, B(H) a C*-algebra GT-bh-cstar), and each JC algebra is a JB algebra (GT-jc-is-jb), so V=B(H)\_sa is a unital JB algebra with Jordan product a∘b=(ab+ba)/2 and order unit 1 (GT-bhsa-jc, GT-jc-is-jb, GT-bh-cstar, GT-jb-orderunit-3310). (a) The polarised hole q\_\{r,s\}:=P(r∘s)-r∘s lies in Ker P: P(q\_\{r,s\})=P(P(r∘s))-P(r∘s)=P\textasciicircum{}2(r∘s)-P(r∘s)=0 since P\textasciicircum{}2=P (GT-Pprops). (b) Fix any state rho on B(H) (rho on V is a state on the order unit space V, GT-jb-orderunit-3310/GT-state-bounded) and set omega:=rho∘Phi. Then omega is a state on V: omega(1)=rho(Phi(1))=rho(1)=1 (Phi unital, GT-positive-unital) and omega(a)=rho(Phi(a))>=0 for a>=0 (Phi positive maps a>=0 to Phi(a)>=0, rho a state, GT-positive-unital). (c) The state seminorm ||x||\_omega\textasciicircum{}2:=omega(x\textasciicircum{}2)=rho(Phi(x\textasciicircum{}2)) is a genuine seminorm on V: omega is a state on the JB algebra V, so a->omega(a\textasciicircum{}2)\textasciicircum{}\{1/2\} is a seminorm (GT-cauchy-schwarz, 3.6.2). In particular it satisfies the triangle inequality and ||lambda x||\_omega=|lambda| ||x||\_omega. Uses GT-Pprops, GT-positive-unital, GT-cauchy-schwarz, GT-bhsa-jc, GT-jc-is-jb, GT-bh-cstar, GT-jb-orderunit-3310, GT-state-bounded, def-near-fixed-algebra. \steptag{validated} \\[4pt]
\step{1.2} (4.1): ||q\_\{r,s\}||\_omega <= C sqrt(eta) ||r|| ||s|| for r,s in A=Im P (def-near-fixed-algebra), every state seminorm omega (node 1.1). PROOF. The map (a,b)->q\_\{a,b\}:=P(a∘b)-a∘b is symmetric (∘ symmetric, def jordan-product) and real-bilinear on A (P real-linear by GT-Pprops, ∘ bilinear; A=Im P is a real subspace so closed under linear combinations); write q\_t:=q\_\{t,t\}=P(t\textasciicircum{}2)-t\textasciicircum{}2 for the diagonal square hole (def-square-hole). Step 1 (polarisation). For a,b in A, a∘b=(1/4)((a+b)\textasciicircum{}2-(a-b)\textasciicircum{}2), so bilinearity+symmetry give q\_\{a,b\}=(1/4)(q\_\{a+b\}-q\_\{a-b\}). For any lambda>0, bilinearity gives q\_\{r,s\}=q\_\{lambda r, s/lambda\}=(1/4)(q\_\{lambda r+s/lambda\}-q\_\{lambda r-s/lambda\}); the arguments lambda r ± s/lambda lie in A (real subspace). Step 2 (seminorm + diagonal bound). By the seminorm triangle inequality (node 1.1) and the DIAGONAL square-hole seminorm bound ||q\_t||\_omega<=C sqrt(eta)||t||\textasciicircum{}2 for t in A (GT-square-hole): ||q\_\{r,s\}||\_omega <= (1/4)(||q\_\{lambda r+s/lambda\}||\_omega+||q\_\{lambda r-s/lambda\}||\_omega) <= (1/4)C sqrt(eta)(||lambda r+s/lambda||\textasciicircum{}2+||lambda r-s/lambda||\textasciicircum{}2). Step 3 (order-unit-norm bound). For any x,y in V, ||x±y||\textasciicircum{}2<=(||x||+||y||)\textasciicircum{}2<=2||x||\textasciicircum{}2+2||y||\textasciicircum{}2 (triangle inequality in the order-unit norm + the scalar bound (p+q)\textasciicircum{}2<=2p\textasciicircum{}2+2q\textasciicircum{}2). With x=lambda r, y=s/lambda: ||lambda r+s/lambda||\textasciicircum{}2+||lambda r-s/lambda||\textasciicircum{}2 <= 4 lambda\textasciicircum{}2||r||\textasciicircum{}2+4 lambda\textasciicircum{}\{-2\}||s||\textasciicircum{}2. Hence ||q\_\{r,s\}||\_omega <= C sqrt(eta)(lambda\textasciicircum{}2||r||\textasciicircum{}2+lambda\textasciicircum{}\{-2\}||s||\textasciicircum{}2). Step 4 (optimise). If r=0 or s=0 then q\_\{r,s\}=0 and the claim is trivial. Otherwise set lambda\textasciicircum{}2=||s||/||r||>0: then lambda\textasciicircum{}2||r||\textasciicircum{}2+lambda\textasciicircum{}\{-2\}||s||\textasciicircum{}2=||r||||s||+||r||||s||=2||r||||s||, giving ||q\_\{r,s\}||\_omega<=2C sqrt(eta)||r||||s||=C' sqrt(eta)||r||||s|| with C'=2C universal dimension-free. Uses node 1.1, GT-square-hole, GT-Pprops, def-near-fixed-algebra, def-square-hole, def jordan-product. \steptag{validated} \\[4pt]
\step{1.3} CRUDE NORM BOUND: ||q\_\{r,s\}|| <= C ||r|| ||s|| (operator/order-unit norm) for r,s in A. PROOF. Here V=B(H)\_sa is a JC algebra (GT-bhsa-jc, B(H) a C*-algebra GT-bh-cstar) and each JC algebra is a JB algebra (GT-jc-is-jb), so V=B(H)\_sa is a unital JB algebra (as set up in node 1.1); this licenses GT-jb-submult below. q\_\{r,s\}=P(r∘s)-r∘s, so by the triangle inequality in the order-unit norm ||q\_\{r,s\}|| <= ||P(r∘s)|| + ||r∘s|| <= ||P|| ||r∘s|| + ||r∘s|| = (||P||+1)||r∘s||. By GT-Pprops ||P||<=1+C eta<=1+1 (eta<=eta\_0<1/4 small, C eta bounded), so ||P||+1<=3 (any universal bound suffices). By the JB-algebra product norm axiom ||r∘s||<=||r|| ||s|| (GT-jb-submult; V=B(H)\_sa is a unital JB algebra, since B(H)\_sa is a JC algebra (GT-bhsa-jc) and each JC algebra is a JB algebra (GT-jc-is-jb), GT-bh-cstar). Hence ||q\_\{r,s\}||<=(||P||+1)||r||||s|| <= (2+C eta)||r||||s|| <= C'||r||||s|| with C'=2+C eta a universal dimension-free constant. Uses GT-Pprops, GT-jb-submult, GT-bhsa-jc, GT-jc-is-jb, GT-bh-cstar, def-near-fixed-algebra. \steptag{validated} \\[6pt]
\step{1.3.1} SEMINORM-DOMINATED-BY-NORM: ||z||\_omega <= ||z|| for every z in V=B(H)\_sa and every state seminorm omega=rho∘Phi (node 1.1). NO use of ||Phi||=1. PROOF. By definition ||z||\_omega\textasciicircum{}2=omega(z\textasciicircum{}2)=rho(Phi(z\textasciicircum{}2)). omega is a state on the JB algebra V (node 1.1), so applying the state bound to the self-adjoint element z\textasciicircum{}2 in V (order unit space V, order unit 1, GT-jb-orderunit-3310): |omega(z\textasciicircum{}2)|<=||z\textasciicircum{}2|| (a state on an order-unit space has |omega(a)|<=||a||, GT-state-bounded). Since z\textasciicircum{}2>=0 and omega is positive, omega(z\textasciicircum{}2)>=0, so omega(z\textasciicircum{}2)=|omega(z\textasciicircum{}2)|<=||z\textasciicircum{}2||. By the JB norm axiom ||z\textasciicircum{}2||=||z||\textasciicircum{}2 (GT-jb-axiom-sq, where V=B(H)\_sa is a JC algebra (GT-bhsa-jc) and each JC algebra is a JB algebra (GT-jc-is-jb), so V is a unital JB algebra). Hence ||z||\_omega\textasciicircum{}2=omega(z\textasciicircum{}2)<=||z||\textasciicircum{}2, and taking square roots ||z||\_omega<=||z||. Uses node 1.1, GT-state-bounded, GT-jb-axiom-sq, GT-jb-orderunit-3310, GT-bhsa-jc, GT-jc-is-jb. \steptag{validated} \\[4pt]
\step{1.4} (HH): ||P(q\_\{r,s\}∘q\_\{u,v\})|| <= C eta ||r||||s||||u||||v|| for r,s,u,v in A. PROOF. Set w:=q\_\{r,s\}∘q\_\{u,v\} in V (q\_\{r,s\},q\_\{u,v\} in V, ∘ closed on V). P(w) is self-adjoint (P:V->V, GT-Pprops). Take the state supremum: ||P(w)||=sup\_\{rho in S(V)\}|rho(P(w))| (GT-sup-states, V order-unit space GT-jb-orderunit-3310). Fix a state rho on V, omega:=rho∘Phi (a state on V, node 1.1). Bridge Phi<->P: |rho(P(w))| <= |rho(Phi(w))| + |rho((P-Phi)(w))| (triangle ineq for scalars). FIRST term: |rho(Phi(w))|=|omega(w)|=|omega(q\_\{r,s\}∘q\_\{u,v\})|. By the Jordan Cauchy-Schwarz inequality for the state omega on the JB algebra V (GT-cauchy-schwarz, 3.6.2(i)): |omega(q\_\{r,s\}∘q\_\{u,v\})| <= omega(q\_\{r,s\}\textasciicircum{}2)\textasciicircum{}\{1/2\} omega(q\_\{u,v\}\textasciicircum{}2)\textasciicircum{}\{1/2\} = ||q\_\{r,s\}||\_omega ||q\_\{u,v\}||\_omega. By (4.1) (node 1.2): ||q\_\{r,s\}||\_omega<=C sqrt(eta)||r||||s|| and ||q\_\{u,v\}||\_omega<=C sqrt(eta)||u||||v||, so |omega(w)| <= C\textasciicircum{}2 eta ||r||||s||||u||||v||. SECOND term (replace Phi by P, the delta cost): |rho((P-Phi)(w))| <= ||(P-Phi)(w)|| (rho a state, |rho(y)|<=||y|| for self-adjoint y=(P-Phi)(w) in V, GT-state-bounded) <= ||P-Phi|| ||w|| = delta ||w|| with delta<=C eta (GT-Pprops). Now ||w||=||q\_\{r,s\}∘q\_\{u,v\}|| <= ||q\_\{r,s\}|| ||q\_\{u,v\}|| (GT-jb-submult) <= (C||r||||s||)(C||u||||v||) (crude bound, node 1.3) = C\textasciicircum{}2||r||||s||||u||||v||. So second term <= C eta · C\textasciicircum{}2 ||r||||s||||u||||v|| = C' eta ||r||||s||||u||||v||. COMBINE: |rho(P(w))| <= (C\textasciicircum{}2+C') eta ||r||||s||||u||||v|| =: C'' eta ||r||||s||||u||||v||, uniformly in rho. Taking the supremum: ||P(q\_\{r,s\}∘q\_\{u,v\})|| <= C'' eta ||r||||s||||u||||v|| with C'' universal dimension-free. Uses node 1.2, node 1.3, node 1.1, GT-cauchy-schwarz, GT-sup-states, GT-state-bounded, GT-jb-submult, GT-Pprops, GT-jb-orderunit-3310. \steptag{validated} \\[4pt]
\step{1.5} (HZ): ||P(q\_\{r,s\}∘z)|| <= C sqrt(eta) ||r||||s||||z|| for r,s in A and z in V=B(H)\_sa. PROOF. Here V=B(H)\_sa is a JC algebra (GT-bhsa-jc) and each JC algebra is a JB algebra (GT-jc-is-jb), so V is a unital JB algebra (as set up in node 1.1); this licenses the JB Cauchy-Schwarz and submultiplicativity used below. Set w:=q\_\{r,s\}∘z in V. P(w) self-adjoint (P:V->V, GT-Pprops), so ||P(w)||=sup\_\{rho in S(V)\}|rho(P(w))| (GT-sup-states, V order-unit space GT-jb-orderunit-3310). Fix a state rho on V, omega:=rho∘Phi (state on V, node 1.1). Bridge: |rho(P(w))| <= |rho(Phi(w))| + |rho((P-Phi)(w))|. FIRST term: |rho(Phi(w))|=|omega(q\_\{r,s\}∘z)| <= ||q\_\{r,s\}||\_omega ||z||\_omega (Jordan Cauchy-Schwarz for the state omega on the JB algebra V, GT-cauchy-schwarz 3.6.2(i)). By (4.1) (node 1.2) ||q\_\{r,s\}||\_omega<=C sqrt(eta)||r||||s||; by node 1.3.1 ||z||\_omega<=||z||. So first term <= C sqrt(eta)||r||||s||·||z|| = C sqrt(eta)||r||||s||||z||. SECOND term: |rho((P-Phi)(w))| <= ||(P-Phi)(w)|| (GT-state-bounded) <= delta||w|| (delta<=C eta, GT-Pprops) and ||w||=||q\_\{r,s\}∘z||<=||q\_\{r,s\}|| ||z|| (GT-jb-submult, V a JB algebra) <= C||r||||s||·||z|| (crude bound node 1.3). So second term <= C eta·C||r||||s||||z|| = C'' eta ||r||||s||||z|| <= C'' sqrt(eta)||r||||s||||z|| (eta<=sqrt(eta) for 0<eta<1). COMBINE: |rho(P(w))| <= (C+C'') sqrt(eta)||r||||s||||z||, uniformly in rho; taking the supremum ||P(q\_\{r,s\}∘z)|| <= C''' sqrt(eta)||r||||s||||z|| with C''' universal dimension-free. Uses node 1.1, node 1.2, node 1.3, node 1.3.1, GT-cauchy-schwarz, GT-sup-states, GT-state-bounded, GT-jb-submult, GT-Pprops, GT-jb-orderunit-3310, GT-bhsa-jc, GT-jc-is-jb. \steptag{validated} \\[4pt]
\end{proofbox}

\end{document}
